\chapter{Process View}
\ChapterLabel{ProcessView}
\section{Chapter Responsibility}
\label{Section:ProcessViewResponsibility}
Explain how the system behaves while it is running.

\noindent\textbf{Expected result:} A runtime model that explains coordination and observable behavior.

\section{Revision 0.2 Scope Notice}
\label{Section:DeferredProcessView}
\begin{fsaNote}[Deferred project content --- template reference only]
This chapter is retained in the complete document as template guidance.
Its sample artifacts, IDs, priorities, statuses, versions, outcomes, and
approval examples do not describe the FTDI245R project. In particular, any
sample Done, Pass, or release entry is not implementation, verification,
validation, or approval evidence for this project. Detailed project content
is deferred under \hyperref[Section:ProjectScope]{the revision scope}.
See \hyperref[Table:ProjectGapReview]{the reviewed gap list} for actual
decisions and unresolved items.
\end{fsaNote}

\section{How to Use This Chapter}
Describe runtime behavior, concurrency, state transitions, workflows, error handling, and significant sequences. Use activity, sequence, or state diagrams when prose is insufficient.
\section{Example State Model}
Request states are Draft, Submitted, Under Review, Approved, Rejected, and Cancelled. Submission moves Draft to Submitted only when required data is valid. Approval and rejection are terminal except for an authorized correction workflow.
